\documentclass[12pt,a4paper]{article}

\usepackage[utf8]{inputenc}
\usepackage[T1]{fontenc}
\usepackage{lmodern}
\usepackage{textcomp}
\usepackage{newunicodechar}
\usepackage[margin=2.5cm]{geometry}
\usepackage{amsmath,amssymb}
\usepackage{graphicx}
\usepackage{listings}
\usepackage{xcolor}
\usepackage{hyperref}
\usepackage{fancyhdr}
\usepackage{float}
\usepackage{booktabs}
\usepackage{caption}

% Adjust header height for fancyhdr
\setlength{\headheight}{15pt}
\addtolength{\topmargin}{-3pt}

% Romanian Unicode support for standard pdflatex
\newunicodechar{Ă}{\v{A}}
\newunicodechar{ă}{\v{a}}
\newunicodechar{Â}{\^{A}}
\newunicodechar{â}{\^{a}}
\newunicodechar{Î}{\^{I}}
\newunicodechar{î}{\^{i}}
\newunicodechar{Ș}{\c{S}}
\newunicodechar{ș}{\c{s}}
\newunicodechar{Ț}{\c{T}}
\newunicodechar{ț}{\c{t}}

\hypersetup{
    colorlinks=true,
    linkcolor=blue,
    filecolor=magenta,      
    urlcolor=cyan,
    citecolor=blue
}

% Code listing configuration matching user preference
\definecolor{codebg}{gray}{0.96}
\definecolor{codegreen}{rgb}{0.13,0.55,0.13}
\definecolor{codegray}{rgb}{0.5,0.5,0.5}
\definecolor{codepurple}{rgb}{0.58,0,0.82}

\lstset{
    language=bash,
    basicstyle=\ttfamily\footnotesize,
    backgroundcolor=\color{codebg},
    frame=single, 
    rulecolor=\color{gray!40},
    breaklines=true, 
    columns=fullflexible,
    keepspaces=true, 
    showstringspaces=false, 
    upquote=true,
    extendedchars=true,
    literate=
        {Ă}{{\v{A}}}1
        {ă}{{\v{a}}}1
        {Â}{{\^{A}}}1
        {â}{{\^{a}}}1
        {Î}{{\^{I}}}1
        {î}{{\^{i}}}1
        {Ș}{{\c{S}}}1
        {ș}{{\c{s}}}1
        {Ț}{{\c{T}}}1
        {ț}{{\c{t}}}1
}

\lstdefinestyle{csharpstyle}{
    language={[Sharp]C},
    backgroundcolor=\color{codebg},   
    commentstyle=\color{codegreen},
    keywordstyle=\color{blue}\bfseries,
    numberstyle=\tiny\color{codegray},
    stringstyle=\color{codepurple},
    basicstyle=\ttfamily\footnotesize,
    breaklines=true,                 
    captionpos=b,                    
    keepspaces=true,                 
    numbers=left,                    
    numbersep=5pt,                  
    showspaces=false,                
    showstringspaces=false,
    tabsize=2,
    frame=single,
    rulecolor=\color{gray!40},
    upquote=true,
    extendedchars=true,
    literate=
        {Ă}{{\v{A}}}1
        {ă}{{\v{a}}}1
        {Â}{{\^{A}}}1
        {â}{{\^{a}}}1
        {Î}{{\^{I}}}1
        {î}{{\^{i}}}1
        {Ș}{{\c{S}}}1
        {ș}{{\c{s}}}1
        {Ț}{{\c{T}}}1
        {ț}{{\c{t}}}1
}

% Header and Footer
\pagestyle{fancy}
\fancyhf{}
\lhead{Cryptography and Security -- Lab 1}
\rhead{Iulian Pavlov, FAF-241}
\cfoot{\thepage}

\begin{document}

% ============================================================================
% 1. TITLE PAGE
% ============================================================================
\begin{titlepage}
\begin{center}
    \includegraphics[width=0.3\textwidth]{logo.png} \\[0.2cm] 
    
    \textbf{MINISTRY OF EDUCATION, CULTURE AND RESEARCH \\ OF THE REPUBLIC OF MOLDOVA} \\[0.3cm]
    
    \textbf{Technical University of Moldova \\ Faculty of Computers, Informatics and Microelectronics \\ Department of Software and Automation Engineering} \\[2cm]
    
    \textbf{Iulian Pavlov, FAF-241}\\[0.5cm]
    
    \Huge \textbf{Report} \\[0.5cm]
    
    \large Laboratory work nr. 1 \\[0.5cm]
    
    \textbf{Classical Encryption Techniques: \\ Caesar Cipher and Caesar Cipher with Permutation} \\[2.8cm]

    \begin{flushright}
        \textit{Checked by:} \\
        \textbf{Zaica Maia}, \textit{University Assistant} \\
        Department of Software and Automation Engineering, \\
        FCIM Faculty, UTM 
    \end{flushright}
    \vfill
    
    Chișinău -- 2026
\end{center}
\end{titlepage}

\newpage
\tableofcontents
\newpage

% ============================================================================
% 2. SHORT THEORY
% ============================================================================
\section{Theoretical Background}

\subsection{The Classical Caesar Cipher}
The Caesar cipher is one of the oldest symmetric cryptosystems, categorized as a \textbf{monoalphabetic substitution cipher}. In this scheme, every letter in the plaintext is shifted by a fixed integer key $k$ across an alphabet $\mathcal{A}$ of cardinality $N = |\mathcal{A}|$.

\subsubsection{Mathematical Formulation}
Let the alphabet symbols be assigned discrete numerical positions in the residue ring $\mathbb{Z}_N = \{0, 1, \dots, N-1\}$.
\begin{itemize}
    \item \textbf{Encryption}:
    \begin{equation}
        E(p, k) = (p + k) \pmod N
    \end{equation}
    where $p \in \mathbb{Z}_N$ is the numerical value of the plaintext letter, $k$ is the integer key, and $N$ is the alphabet length.
    
    \item \textbf{Decryption}:
    \begin{equation}
        D(c, k) = (c - k + N) \pmod N
    \end{equation}
    where $c \in \mathbb{Z}_N$ represents the ciphertext character index.
\end{itemize}

\subsubsection{Key Space Analysis}
For the Romanian alphabet containing $N = 31$ characters (\texttt{A, Ă, Â, B, C, D, E, F, G, H, I, Î, J, K, L, M, N, O, P, Q, R, S, Ș, T, Ț, U, V, W, X, Y, Z}):
\begin{itemize}
    \item The valid non-trivial key space is constrained by:
    \begin{equation}
        \mathcal{K} = \{ k \in \mathbb{Z} \mid 0 < k < 30 \} \implies |\mathcal{K}| = 29
    \end{equation}
    \item Because $|\mathcal{K}| = 29$ is extraordinarily small, an attacker can exhaust all candidate keys almost instantaneously via a brute-force ciphertext-only attack.
\end{itemize}

\subsection{Caesar Cipher with Permutation (Keyed Caesar Cipher)}
To strengthen the classical Caesar cipher against simple inspection, a second secret key---a keyword $k_2$---is utilized to permute the alphabet before applying a numerical shift $k_1$.

\subsubsection{Alphabet Permutation Construction}
Given the standard alphabet $\mathcal{A}$ and a secret keyword $k_2$ (with $|k_2| \ge 7$ containing only alphabet letters):
\begin{enumerate}
    \item Convert $k_2$ to uppercase.
    \item Remove duplicate characters within $k_2$, preserving only their first occurrence.
    \item Append all remaining letters of the original alphabet $\mathcal{A}$ in their original order.
\end{enumerate}
Formally, the permuted alphabet $\mathcal{A}'$ is:
\begin{equation}
    \mathcal{A}' = \operatorname{Distinct}(k_2 \mathbin{\Vert} \mathcal{A})
\end{equation}
where $\mathbin{\Vert}$ denotes concatenation.

\subsubsection{Mathematical Formulation}
Let $\operatorname{index}_{\mathcal{A}'}(x)$ denote the 0-based index of character $x$ in $\mathcal{A}'$:
\begin{itemize}
    \item \textbf{Encryption}:
    \begin{equation}
        c = \mathcal{A}'\left[ (\operatorname{index}_{\mathcal{A}'}(p) + k_1) \pmod N \right]
    \end{equation}
    \item \textbf{Decryption}:
    \begin{equation}
        p = \mathcal{A}'\left[ (\operatorname{index}_{\mathcal{A}'}(c) - k_1 + N) \pmod N \right]
    \end{equation}
\end{itemize}

\subsubsection{Key Space Analysis}
The key consists of the pair $(k_1, k_2)$:
\begin{itemize}
    \item Shift key $k_1 \in \{1, 2, \dots, 29\}$ ($0 < k_1 < 30$).
    \item Keyword $k_2$: With length $|k_2| \ge 7$ chosen from $N = 31$ alphabet characters without repetition, the number of distinct permutations induced solely by keyword prefixes of length $L = 7$ is:
    \begin{equation}
        P(31, 7) = \frac{31!}{(31 - 7)!} = 31 \times 30 \times 29 \times 28 \times 27 \times 26 \times 25 \approx 1.32 \times 10^{10}
    \end{equation}
    Factoring in $k_1$, the composite key space expands to:
    \begin{equation}
        |\mathcal{K}_{\text{total}}| \approx 29 \times 1.32 \times 10^{10} \approx 3.84 \times 10^{11}
    \end{equation}
    This eliminates trivial brute-force recovery. However, being a monoalphabetic substitution, it remains susceptible to frequency analysis.
\end{itemize}

% ============================================================================
% 3. IMPLEMENTATION
% ============================================================================
\section{Implementation Details}

The laboratory tasks were implemented in C\# (.NET 10) organized into modular classes: alphabet definitions, rigorous input validation and logging, mathematical cipher transformations, and task execution pipelines.

\subsection{Alphabet Definition (\texttt{Alphabet.cs})}
The Romanian alphabet consists of 31 characters, accommodating diacritics (\texttt{Ă, Â, Î, Ș, Ț}). Both Romanian and English alphabets are stored as static arrays.

\begin{lstlisting}[style=csharpstyle, caption={Romanian and English alphabets defined in Alphabet.cs}]
namespace Lab1;

public static class Alphabet
{
    public static readonly char[] RomanianAlphabet = 
    {
        'A', 'Ă', 'Â', 'B', 'C', 'D', 'E', 'F', 'G', 'H', 'I', 'Î', 'J', 'K', 'L', 'M',
        'N', 'O', 'P', 'Q', 'R', 'S', 'Ș', 'Ț', 'U', 'V', 'W', 'X', 'Y', 'Z'
    };
    public static readonly char[] EnglishAlphabet = 
    {
        'A', 'B', 'C', 'D', 'E', 'F', 'G', 'H', 'I', 'J', 'K', 'L', 'M',
        'N', 'O', 'P', 'Q', 'R', 'S', 'T', 'U', 'V', 'W', 'X', 'Y', 'Z'
    };
}
\end{lstlisting}
\textit{Comment:} Storing explicit character arrays guarantees deterministic indexing and consistent uppercase manipulation across platforms.

\subsection{Input Validation and Logger (\texttt{logger.cs})}
To fulfill the requirements ($0 < k_1 < 30$, text composed strictly of alphabet characters, and keyword $k_2$ having at least 7 alphabet characters), a centralized input validator was built.

\begin{lstlisting}[style=csharpstyle, caption={Input validation routines in logger.cs}]
namespace Lab1;

public static class Logger
{
    public static bool IsValidKey(int key) => key > 0 && key < 30;

    public static bool IsLetterInAlphabet(char c, char[] alphabet) =>
        Array.IndexOf(alphabet, char.ToUpper(c)) != -1;

    public static bool ValidateText(string text, char[] alphabet, out char invalidChar)
    {
        invalidChar = '\0';
        if (string.IsNullOrEmpty(text)) {
            Console.WriteLine("[Error] Text cannot be empty.");
            return false;
        }
        foreach (char c in text) {
            if (!IsLetterInAlphabet(c, alphabet)) {
                invalidChar = c;
                Console.WriteLine($"[Error] '{c}' is not a valid letter in the alphabet.");
                return false;
            }
        }
        return true;
    }

    public static bool ValidateKey2(string key2, char[] alphabet)
    {
        if (string.IsNullOrEmpty(key2) || key2.Length < 7) {
            Console.WriteLine("[Error] Key 2 must be at least 7 characters long.");
            return false;
        }
        foreach (char c in key2) {
            if (!IsLetterInAlphabet(c, alphabet)) {
                Console.WriteLine($"[Error] Key 2 contains '{c}', which is not a valid letter.");
                return false;
            }
        }
        return true;
    }

    public static int ReadKey(string prompt = "Enter the key (0 < key < 30): ")
    {
        while (true) {
            Console.Write(prompt);
            if (!int.TryParse(Console.ReadLine(), out int key) || !IsValidKey(key)) {
                Console.WriteLine("[Error] Key must satisfy 0 < key < 30 (between 1 and 29).");
                continue;
            }
            return key;
        }
    }
}
\end{lstlisting}
\textit{Comment:} Interactive prompt loops provide clear error feedback upon invalid entry without program termination.

\subsection{Encryption and Decryption Transformations (\texttt{CaesarCypher.cs})}
Modular arithmetic implements shifts in both directions, compensating for negative modulo results.

\begin{lstlisting}[style=csharpstyle, caption={Encryption and Decryption functions in CaesarCypher.cs}]
namespace Lab1;

public static class CaesarCypher
{
    public static char Encryption(char c, int key, char[] alphabet)
    {
        int index = Array.IndexOf(alphabet, c);
        if (index == -1) return c;

        int shift = (index + key) % alphabet.Length;
        if (shift < 0) shift += alphabet.Length;

        return alphabet[shift];
    }

    public static char Decryption(char c, int key, char[] alphabet)
    {
        int index = Array.IndexOf(alphabet, c);
        if (index == -1) return c;

        int shift = (index - key) % alphabet.Length;
        if (shift < 0) shift += alphabet.Length;

        return alphabet[shift];
    }
}
\end{lstlisting}
\textit{Comment:} By receiving the \texttt{alphabet} array as an argument, the exact same implementation serves both Task 1 (standard alphabet) and Task 2 (permuted alphabet).

\subsection{Building the Permuted Alphabet (\texttt{task2.cs})}
Task 2 dynamically derives the permuted alphabet using LINQ \texttt{Concat} and \texttt{Distinct}.

\begin{lstlisting}[style=csharpstyle, caption={Alphabet permutation logic in task2.cs}]
char[] permutedAlphabet = key2.Concat(Alphabet.RomanianAlphabet)
                              .Distinct()
                              .ToArray();
\end{lstlisting}
\textit{Comment:} \texttt{Distinct()} retains only the first occurrence of each character in $k_2 \mathbin{\Vert} \mathcal{A}$, generating a valid bijective permutation of the alphabet.

% ============================================================================
% 4. RESULTS AND VERIFICATION
% ============================================================================
\section{Results and Verification}

To rigorously test and verify the implemented algorithms, comprehensive test cases were selected, reflecting authentic examples from classical cryptography curricula. Both the mathematical step-by-step transformations and the actual program executions are presented below.

\subsection{Task 1: Standard Caesar Cipher}

\subsubsection{Test Case Specification and Mathematical Verification}
\begin{itemize}
    \item \textbf{Alphabet}: Romanian Alphabet ($N = 31$ characters)
    \item \textbf{Secret Key}: $k = 5$ (satisfies $0 < k < 30$)
    \item \textbf{Plaintext}: \texttt{SECURITATE}
\end{itemize}

\begin{table}[H]
\centering
\caption{Mathematical step-by-step verification of Task 1 encryption and decryption}
\label{tab:task1_math}
\resizebox{\textwidth}{!}{%
\begin{tabular}{cccccc}
\toprule
\textbf{Plaintext ($p$)} & \textbf{Index ($p$)} & \textbf{$(p + 5) \pmod{31}$} & \textbf{Ciphertext ($c$)} & \textbf{$(c - 5) \pmod{31}$} & \textbf{Decrypted} \\
\midrule
S & 21 & $(21 + 5) = 26$ & W & $(26 - 5) = 21$ & S \\
E & 6  & $(6 + 5) = 11$  & Î & $(11 - 5) = 6$  & E \\
C & 4  & $(4 + 5) = 9$   & H & $(9 - 5) = 4$   & C \\
U & 25 & $(25 + 5) = 30$ & Z & $(30 - 5) = 25$ & U \\
R & 20 & $(20 + 5) = 25$ & U & $(25 - 5) = 20$ & R \\
I & 10 & $(10 + 5) = 15$ & M & $(15 - 5) = 10$ & I \\
T & 23 & $(23 + 5) = 28$ & X & $(28 - 5) = 23$ & T \\
A & 0  & $(0 + 5) = 5$   & D & $(5 - 5) = 0$   & A \\
T & 23 & $(23 + 5) = 28$ & X & $(28 - 5) = 23$ & T \\
E & 6  & $(6 + 5) = 11$  & Î & $(11 - 5) = 6$  & E \\
\bottomrule
\end{tabular}%
}
\end{table}

\noindent
\textbf{Theoretical Result:} Plaintext \texttt{SECURITATE} maps to ciphertext \texttt{WÎHZUMXDXÎ}, and deciphers precisely back to \texttt{SECURITATE}.

\subsubsection{Program Console Executions}
The program was executed for both encryption and decryption workflows:

\begin{lstlisting}[caption={Task 1 Encryption Execution}]
Enter the key (0 < key < 30): 5
Enter the text: SECURITATE
1. Encrypt text
2. Decrypt text
Enter your choice: 1
Result: WÎHZUMXDXÎ
\end{lstlisting}

\begin{lstlisting}[caption={Task 1 Decryption Execution}]
Enter the key (0 < key < 30): 5
Enter the text: WÎHZUMXDXÎ
1. Encrypt text
2. Decrypt text
Enter your choice: 2
Result: SECURITATE
\end{lstlisting}

\subsection{Task 2: Caesar Cipher with Permutation}

\subsubsection{Test Case Specification and Alphabet Permutation}
\begin{itemize}
    \item \textbf{Shift key ($k_1$)}: $7$ (satisfies $0 < k_1 < 30$)
    \item \textbf{Keyword ($k_2$)}: \texttt{ALGORITM} (length $8 \ge 7$, consisting only of alphabet characters)
    \item \textbf{Plaintext}: \texttt{INFORMATICA}
\end{itemize}

\noindent
\textbf{Permuted Alphabet Construction ($\mathcal{A}'$):}
\begin{enumerate}
    \item Unique keyword characters: \texttt{A, L, G, O, R, I, T, M} ($8$ characters).
    \item Remaining alphabet characters appended in order: \\
    \texttt{Ă, Â, B, C, D, E, F, H, Î, J, K, N, P, Q, S, Ș, Ț, U, V, W, X, Y, Z} ($23$ characters).
    \item Total cardinality: $8 + 23 = 31$ characters.
\end{enumerate}

\begin{table}[H]
\centering
\caption{Mathematical step-by-step verification of Task 2 encryption and decryption}
\label{tab:task2_math}
\resizebox{\textwidth}{!}{%
\begin{tabular}{cccccc}
\toprule
\textbf{Plaintext ($p$)} & \textbf{Index in $\mathcal{A}'$} & \textbf{$(i + 7) \pmod{31}$} & \textbf{Ciphertext ($c$)} & \textbf{$(i_c - 7) \pmod{31}$} & \textbf{Decrypted} \\
\midrule
I & 5  & $(5 + 7) = 12$  & D & $(12 - 7) = 5$  & I \\
N & 19 & $(19 + 7) = 26$ & V & $(26 - 7) = 19$ & N \\
F & 14 & $(14 + 7) = 21$ & Q & $(21 - 7) = 14$ & F \\
O & 3  & $(3 + 7) = 10$  & B & $(10 - 7) = 3$  & O \\
R & 4  & $(4 + 7) = 11$  & C & $(11 - 7) = 4$  & R \\
M & 7  & $(7 + 7) = 14$  & F & $(14 - 7) = 7$  & M \\
A & 0  & $(0 + 7) = 7$   & M & $(7 - 7) = 0$   & A \\
T & 6  & $(6 + 7) = 13$  & E & $(13 - 7) = 6$  & T \\
I & 5  & $(5 + 7) = 12$  & D & $(12 - 7) = 5$  & I \\
C & 11 & $(11 + 7) = 18$ & K & $(18 - 7) = 11$ & C \\
A & 0  & $(0 + 7) = 7$   & M & $(7 - 7) = 0$   & A \\
\bottomrule
\end{tabular}%
}
\end{table}

\noindent
\textbf{Theoretical Result:} Plaintext \texttt{INFORMATICA} encrypts to \texttt{DVQBCFMEDKM} and deciphers symmetrically back to \texttt{INFORMATICA}.

\subsubsection{Program Console Executions}

\begin{lstlisting}[caption={Task 2 Encryption Execution}]
Enter the key (0 < key < 30): 7
Enter the text: INFORMATICA
Enter the second key (at least 7 letters): ALGORITM
Permuted alphabet: A, L, G, O, R, I, T, M, Ă, Â, B, C, D, E, F, H, Î, J, K, N, P, Q, S, Ș, Ț, U, V, W, X, Y, Z
1. Encrypt text
2. Decrypt text
Enter your choice: 1
Result: DVQBCFMEDKM
\end{lstlisting}

\begin{lstlisting}[caption={Task 2 Decryption Execution}]
Enter the key (0 < key < 30): 7
Enter the text: DVQBCFMEDKM
Enter the second key (at least 7 letters): ALGORITM
Permuted alphabet: A, L, G, O, R, I, T, M, Ă, Â, B, C, D, E, F, H, Î, J, K, N, P, Q, S, Ș, Ț, U, V, W, X, Y, Z
1. Encrypt text
2. Decrypt text
Enter your choice: 2
Result: INFORMATICA
\end{lstlisting}

\subsection{Input Validation and Error Handling Tests}
The implementation enforces the following validation checks to guarantee system reliability:
\begin{enumerate}
    \item \textbf{Numerical key bounds ($0 < k < 30$)}: Negative numbers, zero, and values $\ge 30$ are rejected.
    \item \textbf{Plaintext alphabet conformance}: Spaces, numeric digits, and punctuation marks are detected and rejected.
    \item \textbf{Keyword length requirement ($|k_2| \ge 7$)}: Keywords with fewer than 7 characters are rejected.
\end{enumerate}

\begin{lstlisting}[caption={Input Validation Demonstration in Console}]
Enter the key (0 < key < 30): -5
[Error] Key must satisfy 0 < key < 30 (between 1 and 29). You entered: -5
Enter the key (0 < key < 30): 30
[Error] Key must satisfy 0 < key < 30 (between 1 and 29). You entered: 30
Enter the key (0 < key < 30): 42
[Error] Key must satisfy 0 < key < 30 (between 1 and 29). You entered: 42
Enter the key (0 < key < 30): 7
Enter the text: HELLO WORLD 2026!
[Error] ' ' is not a valid letter in the alphabet.
Enter the text: INFORMATICA
Enter the second key (at least 7 letters): CIPHER
[Error] Key 2 must be at least 7 characters long.
Enter the second key (at least 7 letters): UTM
[Error] Key 2 must be at least 7 characters long.
Enter the second key (at least 7 letters): ALGORITM
Permuted alphabet: A, L, G, O, R, I, T, M, Ă, Â, B, C, D, E, F, H, Î, J, K, N, P, Q, S, Ș, Ț, U, V, W, X, Y, Z
1. Encrypt text
2. Decrypt text
Enter your choice: 1
Result: DVQBCFMEDKM
\end{lstlisting}

% ============================================================================
% 5. CONCLUSIONS
% ============================================================================
\section{Conclusions}

In this laboratory work, the classical Caesar cipher and the keyed Caesar cipher with alphabet permutation were studied, implemented, and verified.

Key takeaways and cryptographic insights:
\begin{enumerate}
    \item \textbf{Insecurity of the Classical Caesar Cipher}: Operating over a 31-character alphabet, the Caesar cipher offers only 29 non-trivial keys ($0 < k < 30$), rendering it trivial to crack via brute-force attack in milliseconds.
    
    \item \textbf{Advantage of Keyword Permutation}: Introducing a keyword $k_2$ ($\ge 7$ characters) scrambles the alphabet and expands the key space to over $3.8 \times 10^{11}$ combinations ($P(31, 7) \times 29$), rendering manual brute-force impossible.
    
    \item \textbf{Monoalphabetic Vulnerability}: Because the relationship between plaintext letters and ciphertext letters remains 1-to-1 throughout the encryption of a message, both ciphers remain susceptible to \textbf{frequency analysis}. Natural language statistics (such as high occurrences of letters \texttt{E, I, A} in Romanian) are preserved in the ciphertext.
    
    \item \textbf{Engineering Lessons}: Safe user input handling with iterative error logging, LINQ-based collection deduplication, and modular arithmetic with boundary corrections were successfully implemented in C\#.
\end{enumerate}

\end{document}
